\newpage
\section*{\S~5 Dienstgrade}
\label{sec:dienstgrade}
\addcontentsline{toc}{section}{\S~5 Dienstgrade}

\begin{enumerate}[label=(\arabic*)]
    \item Die 1st Aircat Division übernimmt die Bezeichnung der Dienstgrade der Bundeswehr, aufgrund der deutschsprachigen Ausrichtung der 1st Aircat Division und dem Fokus auf Semi-Milsim.    
    \item Die Dienstgrade dienen dazu, neuen Mitgliedern einen Überblick über Erfahrung, Anwesenheit und Fähigkeiten eines Mitgliedes zu geben.
    \item Die Dienstgrade nehmen, wenn nicht ausdrücklich anders festgelegt, keinen Einfluss auf die Funktion innerhalb einer Mission.
    \item Die Dienstgrade sind durch die angefügte Abkürzung im Arma 3 Namen kenntlich zu machen.  
    \item Die Dienstgrade werden im Folgenden aufgeführt:
    \begin{enumerate}[label=(\Alph*)]
        \item Offiziere
        \begin{itemize}
            \item Hauptmann (Hptm.)
            \item Oberleutnant (OLt.)
            \item Leutnant (Lt.)
        \end{itemize}    
        \item Unteroffiziere mit Portepee (Feldwebel)
        \begin{itemize}
            \item Oberstabsfeldwebel (OStFw.)
            \item Stabsfeldwebel (StFw.)
            \item Hauptfeldwebel (HptFw.)
            \item Oberfeldwebel (OFw.)
            \item Feldwebel (Fw.)
        \end{itemize}  
        \item Unteroffiziere ohne Portepee
        \begin{itemize}
            \item Stabsunteroffizier (StUffz.)
            \item Unteroffizier (Uffz.)
        \end{itemize}
        \item Mannschaften
        \begin{itemize}
            \item Stabskorporal (StKorp.)
            \item Korporal (Korp.)
            \item Oberstabsgefreiter (OStGefr.)
            \item Stabsgefreiter (StGefr.)
            \item Hauptgefreiter (HptGefr.)
            \item Obergefreiter (OGefr.)
            \item Gefreiter (Gefr.)
            \item Jäger (Jg.)
        \end{itemize}
    \end{enumerate}
\end{enumerate}